% Generated draft: requires mathematical and source review.
\documentclass[11pt]{article}
\usepackage[margin=25mm]{geometry}
\usepackage{fontspec}
\setmainfont{Latin Modern Roman}
\usepackage{amsmath,amssymb,mathtools,unicode-math}
\setmathfont{Latin Modern Math}
\usepackage{enumitem,needspace,xurl,hyperref}
\setlength{\parindent}{0pt}
\setlength{\parskip}{5pt}
\setlength{\emergencystretch}{3em}
\newcommand{\N}{\mathbb{N}}
\newcommand{\Z}{\mathbb{Z}}
\newcommand{\Q}{\mathbb{Q}}
\newcommand{\R}{\mathbb{R}}
\newcommand{\C}{\mathbb{C}}
\newcommand{\F}{\mathbb{F}}
\newcommand{\A}{\mathbb{A}}
\newcommand{\PP}{\mathbb{P}}
\newcommand{\E}{\mathbb{E}}
\newcommand{\eps}{\varepsilon}
\newcommand{\dd}{\,\mathrm{d}}
\begin{document}
\section*{UnsolvedMath ID: 30005847}
\textbf{Source problem code:} catalogue entry
\subsection{Definitions and mathematical statement}
Let \(k\) be a coefficient field, placed in degree zero. The supplied record does not specify the coefficient field, the exact structure retained by the pointed graded \(K_0\)-group, or whether classification means graded ring, graded \(k\)-algebra, or graded \(*\)-algebra isomorphism. The formulation below makes these choices explicit as a conventional candidate requiring review; these choices are not additional facts verified from the cited sources.

A finite directed graph is a tuple \(E=(E^0,E^1,s,r)\), where \(E^0\) and \(E^1\) are finite sets of vertices and edges and \(s,r:E^1\to E^0\) give the source and range. A sink is a vertex emitting no edges. Its Leavitt path algebra \(L_k(E)\) is the \(k\)-algebra generated by vertices \(v\in E^0\), edges \(e\in E^1\), and formal ghost edges \(e^*\), subject to
\[
uv=\delta_{u,v}u,\qquad s(e)e=e=e\,r(e),\qquad r(e)e^*=e^*=e^*s(e),
\]
\[
e^*f=\delta_{e,f}r(e),\qquad
v=\sum_{\substack{e\in E^1\\s(e)=v}}ee^*\quad\text{for every vertex }v\text{ that is not a sink}.
\]
In the first relation \(u,v\in E^0\), and in the edge relations \(e,f\in E^1\). The last relation is not imposed at sinks. Its natural \(\Z\)-grading is determined by
\[
\deg(v)=0,\qquad \deg(e)=1,\qquad \deg(e^*)=-1.
\]
For finite \(E\), the identity is \(1=\sum_{v\in E^0}v\). The symbol \(e^*\) here denotes a formal generator; an involution on the coefficient field has not been assumed.

For a unital \(\Z\)-graded algebra \(A\), let \(K_0^{\mathrm{gr}}(A)\) be the Grothendieck group of finitely generated graded projective right \(A\)-modules, using degree-preserving module isomorphisms. Its positive cone \(K_0^{\mathrm{gr}}(A)^+\) consists of classes represented by such modules. Set
\[
P(1)_n=P_{n+1},\qquad x[P]=[P(1)].
\]
Grading shifts give \(K_0^{\mathrm{gr}}(A)\) the structure of a \(\Z[x,x^{-1}]\)-module. Its distinguished element is the class \([A]\) of the regular graded right module. Under the ordered-module interpretation, the pointed invariant is
\[
\mathcal I(A)=\bigl(K_0^{\mathrm{gr}}(A),K_0^{\mathrm{gr}}(A)^+,[A]\bigr),
\]
with the shift action retained. An isomorphism of these invariants means a \(\Z[x,x^{-1}]\)-module isomorphism \(\phi\) satisfying
\[
\phi\bigl(K_0^{\mathrm{gr}}(A)^+\bigr)=K_0^{\mathrm{gr}}(B)^+,
\qquad \phi([A])=[B].
\]

\textbf{Graded Classification Conjecture, conditional ordered-module formulation.} For a fixed coefficient field \(k\), and for every pair of finite directed graphs \(E,F\), is it true that
\[
\mathcal I\bigl(L_k(E)\bigr)\cong\mathcal I\bigl(L_k(F)\bigr)
\quad\Longleftrightarrow\quad
L_k(E)\cong_{\mathrm{gr},k}L_k(F)?
\]
Here the right-hand side means that a unital \(k\)-algebra isomorphism \(\Psi\) exists with
\[
\Psi\bigl(L_k(E)_n\bigr)=L_k(F)_n\qquad(n\in\Z).
\]
This is an existence statement about isomorphism classes. It does not additionally require every prescribed invariant isomorphism \(\phi\) to be induced by a chosen algebra isomorphism. Whether the source intends that stronger lifting requirement is unresolved. Likewise, the source does not explicitly establish that the positive cone and Laurent-polynomial action are part of its phrase ``pointed graded \(K_0\)-group,'' or specify whether the assertion is intended for every field.

\textbf{Supplied scope and partial-result qualifications.} The report passage states that the conjecture holds in special cases, including finite graphs in which every cycle has no exits. An exit from a cycle is an edge emitted by a vertex on that cycle other than the edge followed by the cycle. The passage further reports a result for graphs with a composition series whose factors are cofinal graphs with either a unique sink or a unique cycle without exits. It calls these composition S-NE graphs, with S denoting sink and NE denoting no exits, and describes the same class as graphs with disjoint cycles, finitely many infinite emitters, sinks and cycles, and every infinite path ending in a cycle. An infinite emitter is a vertex emitting infinitely many edges. The reported main result additionally requires each infinite emitter to emit only countably many edges. The meanings of the graph composition series, its factors, and ``ends in a cycle'' are not formally specified in the supplied record, so this description is retained as attributed context rather than used as a new precise theorem.

The passage says that this result includes finite graphs with disjoint cycles, giving the Toeplitz graph as an example, and is also formulated for graph \(C^*\)-algebras. It reports the Isomorphism Conjecture as a consequence for the class considered, an equivalent graph-theoretic condition for the S-NE graphs considered, and an S-NE-invariant describing their graded \((*\!-)\)isomorphism classes. The precise statements and hypotheses of those related assertions are not supplied. It also suggests induction on composition-series length as a possible approach to other graph classes. These contextual assertions do not restrict the finite-graph target to the classes already treated.

The supplied dated assessment labels the problem partially solved, attributes results to disjoint-cycle and composition S-NE classes, mentions further primitive-graph homotopy results, and says the full finite-graph conjecture remains open. This is supplied curation, not an independently verified mathematical or literature-status conclusion.
\subsection{Short English statement}
The question asks whether a naturally graded Leavitt path algebra of a finite graph is determined, up to graded isomorphism, by its pointed graded \(K_0\)-invariant. A conventional precise version retains the positive cone, grading-shift action, and regular-module class over a fixed coefficient field. The supplied record leaves those structural choices and the intended type of isomorphism unresolved. It reports partial results for special graph classes; its general open-status assessment requires independent review.
\subsection{Sources}
\begin{itemize}
\item Supplied local UnsolvedMath registry record, permanent ID \(30005847\), title ``Graded Classification Conjecture for Leavitt Path Algebras,'' exact source problem code \texttt{catalogue entry}, in the Oberwolfach Reports Open Problems collection. The record gives proposed year 2024 and supplies no named proposer.
\item Source report: \emph{Combinatorial \(*\)-algebras} (2024), Oberwolfach Reports, pp. 731--788, DOI 10.4171/owr/2024/12. Supplied link: \url{https://doi.org/10.4171/owr/2024/12}. The supplied original statement is a report passage retained as context; no individual contribution author or precise problem page is identified.
\item Supplied literature source: \emph{Graded homotopy classification of Leavitt path algebras}. \url{https://doi.org/10.4171/RMI/1515}.
\item Supplied literature source: \emph{The graded classification conjecture for disjoint-cycle graphs}. \url{https://liavas.net/files/Disjoint_cycles.pdf}.
\item Literature assessment dated exactly \texttt{2026-08-22}: supplied dated curation, not independently verified. It labels the problem partially solved, reports results for disjoint-cycle and composition S-NE graphs and further primitive-graph homotopy results, and states that the full finite-graph conjecture remains open. The supplied background explicitly requests independent verification of time-sensitive status and cited claims.
\item The registry labels the statement \texttt{corrected\_verified}, explaining only that formula delimiters, notation, grammar, and typography were repaired while mathematical content and scope were preserved. This supplied label does not establish mathematical correctness or independent source verification of this draft.
\end{itemize}
\end{document}
